
\documentclass[11pt,a4paper]{article}
\usepackage[utf8]{inputenc}
\usepackage[T1]{fontenc}
\usepackage{amsmath, amssymb, amsfonts}
\usepackage{geometry}
\geometry{margin=1in}

\begin{document}

\section{Computational Examples}

\subsection{Example 1 --- Real Distinct Eigenvalues}
\textbf{System:} $\dot{x}(t) = Ax(t)$, where $A = \begin{bmatrix} 1 & 2 \\ 2 & 1 \end{bmatrix}$.

\paragraph{Step 1 --- Eigenvalues}
The characteristic equation is given by:
\begin{equation*}
\det(A - \lambda I) = (1-\lambda)^2 - 4 = \lambda^2 - 2\lambda - 3 = (\lambda-3)(\lambda+1) = 0
\end{equation*}
Thus, the eigenvalues are: $\lambda_1 = 3, \lambda_2 = -1$.

\paragraph{Step 2 --- Eigenvectors}
\begin{itemize}
    \item For $\lambda_1 = 3$: Solve $(A-3I)v = 0 \implies \begin{bmatrix} -2 & 2 \\ 2 & -2 \end{bmatrix} \begin{bmatrix} v_{11} \\ v_{12} \end{bmatrix} = 0 \implies v_1 = \begin{bmatrix} 1 \\ 1 \end{bmatrix}$.
    \item For $\lambda_2 = -1$: Solve $(A+I)v = 0 \implies \begin{bmatrix} 2 & 2 \\ 2 & 2 \end{bmatrix} \begin{bmatrix} v_{21} \\ v_{22} \end{bmatrix} = 0 \implies v_2 = \begin{bmatrix} 1 \\ -1 \end{bmatrix}$.
\end{itemize}

\paragraph{Step 3 --- Matrices $P$ and $D$}
\begin{equation*}
P = \begin{bmatrix} 1 & 1 \\ 1 & -1 \end{bmatrix}, \quad D = \begin{bmatrix} 3 & 0 \\ 0 & -1 \end{bmatrix}
\end{equation*}

\paragraph{Step 4 --- Invert $P$}
\begin{equation*}
P^{-1} = \frac{1}{-2} \begin{bmatrix} -1 & -1 \\ -1 & 1 \end{bmatrix} = \frac{1}{2} \begin{bmatrix} 1 & 1 \\ 1 & -1 \end{bmatrix}
\end{equation*}

\paragraph{Step 5 \& 6 --- Assemble $e^{At}$}
\begin{equation*}
e^{Dt} = \begin{bmatrix} e^{3t} & 0 \\ 0 & e^{-t} \end{bmatrix}
\end{equation*}
\begin{equation*}
e^{At} = P e^{Dt} P^{-1} = \frac{1}{2} \begin{bmatrix} e^{3t}+e^{-t} & e^{3t}-e^{-t} \\ e^{3t}-e^{-t} & e^{3t}+e^{-t} \end{bmatrix}
\end{equation*}
\textbf{Verification:} At $t = 0$: $e^{A \cdot 0} = \frac{1}{2}\begin{bmatrix} 2 & 0 \\ 0 & 2 \end{bmatrix} = I$ \checkmark

\subsection{Example 2 --- Pure Imaginary Eigenvalues}
\textbf{System:} $\dot{x}(t) = Ax(t)$, where $A = \begin{bmatrix} 0 & -1 \\ 1 & 0 \end{bmatrix}$ (rotation matrix).

\paragraph{Step 1 --- Eigenvalues}
\begin{equation*}
\det(A - \lambda I) = \lambda^2 + 1 = 0 \implies \lambda_1 = i, \lambda_2 = -i
\end{equation*}

\paragraph{Steps 2--6 --- Compute $e^{At}$ using Euler's formula}
Applying the full procedure yields:
\begin{equation*}
e^{At} = \begin{bmatrix} \cos t & -\sin t \\ \sin t & \cos t \end{bmatrix}
\end{equation*}
\textbf{Interpretation:} $e^{At}$ is the standard 2D rotation matrix by angle $t$. The purely imaginary eigenvalues $\lambda = \pm i$ (i.e., $\alpha = 0, \beta = 1$) produce sustained oscillation with no growth or decay — a marginally stable system.

\textbf{Verification:} At $t = 0$: $e^{A \cdot 0} = \begin{bmatrix} 1 & 0 \\ 0 & 1 \end{bmatrix} = I$ \checkmark

\section{Role of $e^{At}$ in the General Solution}
The matrix exponential $e^{At}$ is the unique fundamental matrix satisfying $e^{A \cdot 0} = I$. It acts as the state-transition operator of the linear system: the complete solution of $\dot{x}(t) = Ax(t) + f(t)$ is given by the \textit{Variation of Parameters} formula:
\begin{equation}
x(t) = e^{At}x(0) + \int_{0}^{t} e^{A(t-s)} f(s) \, ds
\end{equation}
The first term $e^{At}x(0)$ is the homogeneous response (free evolution from initial conditions); the integral term is the particular response (forced evolution due to $f(t)$). The convolution structure shows that the current state depends on the entire history of the input, filtered through $e^{A(t-s)}$.

\section{References}
\begin{enumerate}
    \item Horn, R. A., \& Johnson, C. R. (2013). \textit{Matrix Analysis} (2nd ed.). Cambridge University Press. [Definition 2.1, Theorem 2.1 — pp. 439–444]
    \item Hirsch, M. W., Smale, S., \& Devaney, R. L. (2004). \textit{Differential Equations, Dynamical Systems, and an Introduction to Chaos} (2nd ed.). Elsevier Academic Press. [Semigroup property — p. 134]
    \item Hall, B. C. (2015). \textit{Lie Groups, Lie Algebras, and Representations: An Elementary Introduction} (2nd ed.). Springer. [Baker–Campbell–Hausdorff formula — Ch. 3]
    \item Strang, G. (2016). \textit{Introduction to Linear Algebra} (5th ed.). Wellesley-Cambridge Press. [Motivation, diagonalization — Ch. 6]
    \item Axler, S. (2015). \textit{Linear Algebra Done Right} (3rd ed.). Springer. [Spectral theorem, diagonalizability conditions — Ch. 8]
\end{enumerate}

\end{document}
